\section{Boundaries}

The claim boundary is the one stated in the V.8 plan. RCOUNT owns count
summaries and reconciliation. RPLAN owns assignments and plan hashes.
\texttt{rplan-core} validates the plan and computes the assignment-bearing plan
hash. RCTX owns context and crosswalk identity; \texttt{rctx-core} verifies the
crosswalk record set and computes its hash. The bridge composes these hashes
and sums counts. It does not make RCOUNT a plan-validity verifier.

The implemented crosswalk path is intentionally conservative. It can use
explicit RCTX crosswalk rows and reject declared hash drift, but the transcript
does not yet record non-exhaustive crosswalk caveats. Rows whose target unit is
not in the plan are ignored by the current aggregation path, so a future public
fixture should make the non-exhaustive boundary visible rather than treating it
as a silent policy choice.

Fractional allocation is also bounded. The bridge accepts rational crosswalk
weights only when every affected count field allocates to an integer. That is
appropriate for a count transcript that emits integer RCOUNT summaries, but it
is not a general fractional redistricting report. A rational-output transcript
mode remains future work for genuinely fractional population, area, or vote
allocation records.

Finally, the implemented fixtures are synthetic. They prove artifact binding,
lineage failure, stale-plan failure, explicit-crosswalk projection, and
crosswalk-hash drift. They do not yet prove a real public precinct-to-district
crosswalk ingestion path.
